\chapter{Study Materials and Descriptive Tables}
\label{app:materials-descriptives}

This appendix groups the reproducible study artifacts (the
standardized scenario and the verbatim items of the user-study
instruments), the full demographic tables, and the per-item
descriptive distributions that complement the results reported
in Chapters~\ref{ch:userstudy} and~\ref{ch:docquality}.

\section{Study Scenario}
\label{app:scenario}

The scenario reproduced below was presented to all participants
in German, regardless of interaction mode (Form, Ask, Chat). The
English translation is provided for the international reader; the
German original is the version actually shown.

\subsection{Original (German)}
\label{app:scenario-de}

\begingroup
\lstset{
  basicstyle=\ttfamily\scriptsize,
  breaklines=true,
  breakatwhitespace=false,
  breakindent=0pt,
  postbreak=\mbox{$\hookrightarrow$\space},
  frame=single,
  framerule=0.4pt,
  language=,
  numbers=none,
  showstringspaces=false,
  keepspaces=true,
  columns=fullflexible,
  upquote=true,
  aboveskip=0.6em,
  belowskip=0.6em,
  literate=%
    {ä}{{\"a}}1 {ö}{{\"o}}1 {ü}{{\"u}}1
    {Ä}{{\"A}}1 {Ö}{{\"O}}1 {Ü}{{\"U}}1 {ß}{{\ss}}1
    {§}{{\S}}1
    {•}{{\textbullet}}1
    {→}{{$\to$}}1
    {≥}{{$\geq$}}1 {≤}{{$\leq$}}1
    {–}{{--}}1 {—}{{---}}1
    {“}{{``}}1 {”}{{''}}1
    {‘}{{`}}1 {’}{{'}}1
}
\begin{lstlisting}
Szenario: Mitarbeiterverwaltung
(Arbeitszeit/Urlaub) & Parkplatzberechtigung

Sie sind Mitarbeitender namens A.I. und arbeiten
für die APOR GmbH, ein Softwareunternehmen mit
ca. 50 Mitarbeitenden in Ulm. Sie handeln im Namen
des Unternehmens und sind im Szenario sowohl
Verantwortlicher (Data Controller) als auch
Datenschutzbeauftragte:r (DPO/DSB).

Zur Organisation:

Rolle: Data Controller, DPO
Adresse: Mainstraße 67, Ulm
E-Mail: mail@apor.eu
Telefon: 012345678
Vertreter: (Sie), A. I., ai@apor.eu
DPO: Das sind auch Sie

Verarbeitung
Das Unternehmen verarbeitet personenbezogene Daten,
um Mitarbeitende vom Eintritt bis zum Austritt zu
verwalten. Im Fokus stehen:

- Anlegen neuer Mitarbeitender
- Dokumentation von Arbeitszeiten
- Verwaltung von Urlaub und Abwesenheiten

Zusätzlich verwaltet APOR für berechtigte
Mitarbeitende eine Parkplatzberechtigung für einen
Firmenparkplatz mit Schranke. Damit berechtigte
Mitarbeitende den Parkplatz nutzen können, wird
ein geeignetes Merkmal zur Zuordnung der
Berechtigung hinterlegt (z. B. ein
Fahrzeugkennzeichen).

APOR arbeitet mit intern gehostetem Microsoft
Office sowie internen Servern (keine externe
Cloud). Krankentage oder medizinische Angaben sind
für diese Verarbeitung nicht erforderlich.

Für die Verwaltung sind technische und
organisatorische Maßnahmen vorgesehen, wie z.B.
Zugriffsbeschränkung nach Rollen (Need-to-know),
Passwortschutz (ggf. MFA), regelmäßige Backups,
Zutrittskontrolle zu Serverraum/Archiv und
Vertraulichkeitsverpflichtung der berechtigten
Mitarbeitenden.

Alle in diesem Zusammenhang gespeicherten
Informationen werden spätestens 3 Monate nach
Austritt der jeweiligen Person gelöscht.
\end{lstlisting}
\endgroup

\subsection{English Translation}
\label{app:scenario-en}

\begingroup
\lstset{
  basicstyle=\ttfamily\scriptsize,
  breaklines=true,
  breakatwhitespace=false,
  breakindent=0pt,
  postbreak=\mbox{$\hookrightarrow$\space},
  frame=single,
  framerule=0.4pt,
  language=,
  numbers=none,
  showstringspaces=false,
  keepspaces=true,
  columns=fullflexible,
  upquote=true,
  aboveskip=0.6em,
  belowskip=0.6em,
  literate=%
    {ä}{{\"a}}1 {ö}{{\"o}}1 {ü}{{\"u}}1
    {Ä}{{\"A}}1 {Ö}{{\"O}}1 {Ü}{{\"U}}1 {ß}{{\ss}}1
    {§}{{\S}}1
    {•}{{\textbullet}}1
    {→}{{$\to$}}1
    {≥}{{$\geq$}}1 {≤}{{$\leq$}}1
    {–}{{--}}1 {—}{{---}}1
    {“}{{``}}1 {”}{{''}}1
    {‘}{{`}}1 {’}{{'}}1
}
\begin{lstlisting}
Scenario: Employee Administration
(Working Time / Leave) & Parking Authorization

You are an employee named A.I., working for APOR
GmbH, a software company with approximately 50
employees in Ulm. You act on behalf of the company
and, in this scenario, are both the Data Controller
and the Data Protection Officer (DPO).

Organization details:

Role: Data Controller, DPO
Address: Mainstrasse 67, Ulm
E-mail: mail@apor.eu
Phone: 012345678
Representative: (You), A. I., ai@apor.eu
DPO: Also you

Processing
The company processes personal data in order to
administer employees from onboarding to exit. The
focus is on:

- Onboarding new employees
- Documentation of working hours
- Administration of leave and absences

In addition, APOR manages a parking authorization
for entitled employees for a company car park with
a barrier. To allow entitled employees to use the
car park, a suitable feature for assigning the
authorization is recorded (e.g. a vehicle license
plate).

APOR uses internally hosted Microsoft Office and
internal servers (no external cloud). Sick days or
medical information are not required for this
processing.

For administration, technical and organizational
measures (TOMs) are provided, such as role-based
access restriction (need-to-know), password
protection (where appropriate, MFA), regular
backups, physical access control to the server
room / archive, and a confidentiality obligation
for authorized employees.

All information stored in this context is deleted
no later than 3 months after the person's
departure.
\end{lstlisting}
\endgroup


\section{User-Study Instruments}
\label{app:sus-rtlx-items}

This section reproduces the verbatim items of the two standardized
user-study instruments referenced from the Methodology chapter:
the \textit{System Usability Scale} (\textit{SUS}) and the \textit{Raw NASA Task Load Index}
(\textit{R-TLX}). The six-item \textit{Perceived AI Support} instrument is bespoke
to this study and is listed in full in the Methodology chapter;
it is not reproduced here.

\subsection{System Usability Scale (SUS)}
\label{app:sus-items}

Five-point Likert scale, $1 = $ Strongly Disagree, $5 = $ Strongly
Agree. Odd-numbered items are positively keyed; even-numbered
items are negatively keyed. The canonical wording reproduced below
is taken from Brooke~\cite{Brooke1996}.

\begin{enumerate}
    \item I think that I would like to use this system frequently.
    \item I found the system unnecessarily complex.
    \item I thought the system was easy to use.
    \item I think that I would need the support of a technical person to be able to use this system.
    \item I found the various functions in this system were well integrated.
    \item I thought there was too much inconsistency in this system.
    \item I would imagine that most people would learn to use this system very quickly.
    \item I found the system very cumbersome to use.
    \item I felt very confident using the system.
    \item I needed to learn a lot of things before I could get going with this system.
\end{enumerate}

\subsection{NASA Raw Task Load Index Subscales}
\label{app:rtlx-items}

Six subscales, each rated on the standard NASA TLX 21-point
scale. The Raw TLX variant used in this study takes the
unweighted mean across the six subscales rather than applying the
pairwise-comparison weighting of the original NASA TLX. The
subscale definitions reproduced below follow Hart~\cite{Hart2006}.

\begin{description}
    \item[Mental Demand] How mentally demanding was the task?
    \item[Physical Demand] How physically demanding was the task?
    \item[Temporal Demand] How hurried or rushed was the pace of the task?
    \item[Performance] How successful were you in accomplishing what you were asked to do?
    \item[Effort] How hard did you have to work to accomplish your level of performance?
    \item[Frustration] How insecure, discouraged, irritated, stressed, and annoyed were you?
\end{description}


\section{Participant Demographics in Full}
\label{app:demographics}

This section reproduces the descriptive demographic tables
referenced in Section~\ref{sec:eval-demographics}. The numeric and
ordinal variables (session duration and three Likert self-reports
of prior expertise) are tabulated first, followed by the
categorical variables (gender, age, degree, profession, field).

\subsection{Numeric and Ordinal Variables}
\label{app:demographics-numeric}

\begin{table}[H]
\centering
\small
\begin{tabular}{@{}l r r r r r r r r@{}}
\toprule
Variable & $n$ & Mean & SD & Q1 & Median & Q3 \\
\midrule
Duration (s)         & 69 & 1865.7 & 1354.2 & 840 & 2013 & 2708 \\
GDPR knowledge       & 69 & 2.81  & 1.15  & 2   & 3    & 4 \\
ROPA familiarity     & 69 & 1.45  & 0.76  & 1   & 1    & 2 \\
Chatbot familiarity  & 69 & 4.36  & 0.73  & 4   & 4    & 5 \\
\bottomrule
\end{tabular}
\caption{Descriptive statistics for numeric and ordinal demographic
variables ($n = 69$). Duration is in seconds and aggregates the full
study session across all three modes; the three expertise variables
are five-point Likert self-reports (1 = no familiarity, 5 = expert).}
\label{tab:demographics-numeric}
\end{table}

\subsection{Categorical Variables}
\label{app:demographics-categorical}

\begin{table}[H]
\centering
\small
\begin{tabular}{@{}l l r r@{}}
\toprule
Variable & Level & $n$ & Share \\
\midrule
\multirow{3}{*}{Gender ($n = 66$ reporters)}
  & Male           & 41 & 62.1\,\% \\
  & Female         & 24 & 36.4\,\% \\
  & Other          &  1 & \phantom{0}1.5\,\% \\
\midrule
\multirow{4}{*}{Age (years)}
  & 24             & 18 & 26.1\,\% \\
  & 25             & 18 & 26.1\,\% \\
  & 26             & 10 & 14.5\,\% \\
  & Other (21--40) & 23 & 33.3\,\% \\
\midrule
\multirow{3}{*}{Degree ($n = 68$ reporters)}
  & Bachelor       & 60 & 88.2\,\% \\
  & Master         &  6 & \phantom{0}8.8\,\% \\
  & Other          &  2 & \phantom{0}2.9\,\% \\
\midrule
\multirow{4}{*}{Profession}
  & Student        & 56 & 81.2\,\% \\
  & Employed       & 10 & 14.5\,\% \\
  & PhD            &  2 & \phantom{0}2.9\,\% \\
  & Staatsexamen   &  1 & \phantom{0}1.5\,\% \\
\midrule
\multirow{5}{*}{Field}
  & IT             & 41 & 59.4\,\% \\
  & Business (BWL) & 16 & 23.2\,\% \\
  & Other          &  6 & \phantom{0}8.7\,\% \\
  & Research       &  4 & \phantom{0}5.8\,\% \\
  & HR             &  2 & \phantom{0}2.9\,\% \\
\bottomrule
\end{tabular}
\caption{Categorical demographic variables ($n = 69$ unless noted).
Shares are computed within variable.}
\label{tab:demographics-categorical}
\end{table}


\section{User-Study Per-Item Descriptive Tables}
\label{app:user-study-items}

This section reproduces the full per-item descriptive
distributions for each user-study instrument. Each per-mode column
reports the mean response on the original Likert scale ($n = 69$
per mode); the Aggregated column averages the three per-mode means.

\subsection{System Usability Scale}
\label{app:items-sus}

The compact per-item means table and the box plots are reproduced
inline with the SUS results in Section~\ref{sec:eval-sus}; the
full per-item descriptive distribution is reproduced here.

\begin{sidewaystable}
\centering
\footnotesize
\begin{tabularx}{\textheight}{@{}c X l r r r r r r r@{}}
\toprule
\# & Item & Mode & Mean & SD & Q1 & Median & Q3 & Min & Max \\
\midrule
\multirow{4}{*}{1}  & \multirow{4}{*}{I think that I would like to use this system frequently.}
                    & Aggregated  & 3.32 & 1.14 & 3 & 3 & 4 & 1 & 5 \\
                    & & Form & 3.52 & 1.09 & 3 & 4 & 4 & 1 & 5 \\
                    & & Ask  & 3.22 & 1.08 & 3 & 3 & 4 & 1 & 5 \\
                    & & Chat & 3.22 & 1.24 & 2 & 3 & 4 & 1 & 5 \\
\midrule
\multirow{4}{*}{2}  & \multirow{4}{*}{I found the system unnecessarily complex.}
                    & Aggregated  & 2.00 & 1.12 & 1 & 2 & 3 & 1 & 5 \\
                    & & Form & 1.64 & 0.92 & 1 & 1 & 2 & 1 & 5 \\
                    & & Ask  & 2.25 & 1.22 & 1 & 2 & 3 & 1 & 5 \\
                    & & Chat & 2.12 & 1.12 & 1 & 2 & 3 & 1 & 5 \\
\midrule
\multirow{4}{*}{3}  & \multirow{4}{*}{I thought the system was easy to use.}
                    & Aggregated  & 4.13 & 0.97 & 4 & 4 & 5 & 1 & 5 \\
                    & & Form & 4.38 & 0.91 & 4 & 5 & 5 & 1 & 5 \\
                    & & Ask  & 4.01 & 0.81 & 4 & 4 & 5 & 1 & 5 \\
                    & & Chat & 3.99 & 1.13 & 3 & 4 & 5 & 1 & 5 \\
\midrule
\multirow{4}{*}{4}  & \multirow{4}{=}{I think that I would need the support of a technical person to be able to use this system.}
                    & Aggregated  & 1.52 & 0.89 & 1 & 1 & 2 & 1 & 5 \\
                    & & Form & 1.52 & 0.85 & 1 & 1 & 2 & 1 & 5 \\
                    & & Ask  & 1.49 & 0.80 & 1 & 1 & 2 & 1 & 5 \\
                    & & Chat & 1.55 & 1.01 & 1 & 1 & 2 & 1 & 5 \\
\midrule
\multirow{4}{*}{5}  & \multirow{4}{=}{I found the various functions in this system were well integrated.}
                    & Aggregated  & 3.77 & 1.05 & 3 & 4 & 5 & 1 & 5 \\
                    & & Form & 4.00 & 0.92 & 4 & 4 & 5 & 1 & 5 \\
                    & & Ask  & 3.67 & 1.02 & 3 & 4 & 4 & 1 & 5 \\
                    & & Chat & 3.64 & 1.16 & 3 & 4 & 5 & 1 & 5 \\
\bottomrule
\end{tabularx}
\caption{System Usability Scale, full per-item descriptive
distribution by mode, items 1--5 ($n = 69$ per mode; see
Table~\ref{tab:sus-per-item-distribution-b} for items 6--10).
Source: \texttt{sus\_per\_item\_summary.csv}.}
\label{tab:sus-per-item-distribution-a}
\end{sidewaystable}
\FloatBarrier

\begin{sidewaystable}
\centering
\footnotesize
\begin{tabularx}{\textheight}{@{}c X l r r r r r r r@{}}
\toprule
\# & Item & Mode & Mean & SD & Q1 & Median & Q3 & Min & Max \\
\midrule
\multirow{4}{*}{6}  & \multirow{4}{*}{I thought there was too much inconsistency in this system.}
                    & Aggregated  & 2.07 & 1.00 & 1 & 2 & 3 & 1 & 5 \\
                    & & Form & 2.16 & 0.96 & 1 & 2 & 3 & 1 & 5 \\
                    & & Ask  & 2.04 & 1.06 & 1 & 2 & 3 & 1 & 5 \\
                    & & Chat & 2.01 & 0.98 & 1 & 2 & 3 & 1 & 5 \\
\midrule
\multirow{4}{*}{7}  & \multirow{4}{=}{I would imagine that most people would learn to use this system very quickly.}
                    & Aggregated  & 4.11 & 0.98 & 4 & 4 & 5 & 1 & 5 \\
                    & & Form & 4.30 & 0.96 & 4 & 5 & 5 & 1 & 5 \\
                    & & Ask  & 4.00 & 0.95 & 4 & 4 & 5 & 1 & 5 \\
                    & & Chat & 4.03 & 1.01 & 4 & 4 & 5 & 1 & 5 \\
\midrule
\multirow{4}{*}{8}  & \multirow{4}{*}{I found the system very cumbersome to use.}
                    & Aggregated  & 2.03 & 1.13 & 1 & 2 & 3 & 1 & 5 \\
                    & & Form & 1.65 & 0.94 & 1 & 1 & 2 & 1 & 5 \\
                    & & Ask  & 2.16 & 1.12 & 1 & 2 & 3 & 1 & 5 \\
                    & & Chat & 2.29 & 1.23 & 1 & 2 & 3 & 1 & 5 \\
\midrule
\multirow{4}{*}{9}  & \multirow{4}{*}{I felt very confident using the system.}
                    & Aggregated  & 3.86 & 1.04 & 3 & 4 & 5 & 1 & 5 \\
                    & & Form & 3.83 & 1.04 & 3 & 4 & 5 & 1 & 5 \\
                    & & Ask  & 3.84 & 0.93 & 3 & 4 & 4 & 1 & 5 \\
                    & & Chat & 3.90 & 1.15 & 4 & 4 & 5 & 1 & 5 \\
\midrule
\multirow{4}{*}{10} & \multirow{4}{=}{I needed to learn a lot of things before I could get going with this system.}
                    & Aggregated  & 1.61 & 0.90 & 1 & 1 & 2 & 1 & 5 \\
                    & & Form & 1.61 & 0.86 & 1 & 1 & 2 & 1 & 5 \\
                    & & Ask  & 1.61 & 0.81 & 1 & 1 & 2 & 1 & 5 \\
                    & & Chat & 1.62 & 1.03 & 1 & 1 & 2 & 1 & 5 \\
\bottomrule
\end{tabularx}
\caption{System Usability Scale, full per-item descriptive
distribution by mode, items 6--10 ($n = 69$ per mode; continued
from Table~\ref{tab:sus-per-item-distribution-a}). Source:
\texttt{sus\_per\_item\_summary.csv}.}
\label{tab:sus-per-item-distribution-b}
\end{sidewaystable}
\FloatBarrier

\begin{figure}[!htbp]
\centering
\includegraphics[width=0.45\linewidth]{sus_item_01.png}\hfill
\includegraphics[width=0.45\linewidth]{sus_item_02.png}\\[1ex]
\includegraphics[width=0.45\linewidth]{sus_item_03.png}\hfill
\includegraphics[width=0.45\linewidth]{sus_item_04.png}\\[1ex]
\includegraphics[width=0.45\linewidth]{sus_item_05.png}\hfill
\includegraphics[width=0.45\linewidth]{sus_item_06.png}
\caption{Box plots of SUS items 1--6 by mode ($n = 69$ per mode),
one plot per item in administered order. Boxes show the
interquartile range, the central line marks the median, and the
diamond marks the mean. Continues in
Figure~\ref{fig:sus-subquestions-b}.}
\label{fig:sus-subquestions-a}
\end{figure}

\begin{figure}[!htbp]
\centering
\includegraphics[width=0.45\linewidth]{sus_item_07.png}\hfill
\includegraphics[width=0.45\linewidth]{sus_item_08.png}\\[1ex]
\includegraphics[width=0.45\linewidth]{sus_item_09.png}\hfill
\includegraphics[width=0.45\linewidth]{sus_item_10.png}
\caption{Box plots of SUS items 7--10 by mode ($n = 69$ per mode),
one plot per item in administered order. Continued from
Figure~\ref{fig:sus-subquestions-a}.}
\label{fig:sus-subquestions-b}
\end{figure}

\subsection{NASA Raw Task Load Index}
\label{app:items-rtlx}

The compact per-subscale means table and the box plots are
reproduced inline with the R-TLX results in
Section~\ref{sec:eval-rtlx}; the full per-subscale descriptive
distribution is reproduced here.

\begin{table}[H]
\centering
\small
\begin{tabular}{@{}l l r r r r r r r@{}}
\toprule
Subscale & Mode & Mean & SD & Q1 & Median & Q3 & Min & Max \\
\midrule
\multirow{4}{*}{Mental Demand}     & Aggregated  & 6.31 & 4.75 & 2.0 & 5.0  & 9.5  & 1 & 21 \\
                                   & Form & 5.77 & 4.26 & 2.0 & 5.0  & 9.0  & 1 & 17 \\
                                   & Ask  & 6.80 & 4.63 & 3.0 & 6.0  & 10.0 & 1 & 21 \\
                                   & Chat & 6.36 & 5.31 & 2.0 & 5.0  & 10.0 & 1 & 21 \\
\midrule
\multirow{4}{*}{Physical Demand}   & Aggregated  & 3.19 & 3.40 & 1.0 & 2.0  & 4.0  & 1 & 21 \\
                                   & Form & 2.77 & 2.71 & 1.0 & 1.0  & 4.0  & 1 & 14 \\
                                   & Ask  & 3.54 & 3.92 & 1.0 & 2.0  & 5.0  & 1 & 21 \\
                                   & Chat & 3.26 & 3.46 & 1.0 & 2.0  & 4.0  & 1 & 20 \\
\midrule
\multirow{4}{*}{Temporal Demand}   & Aggregated  & 4.89 & 4.28 & 1.0 & 4.0  & 7.0  & 1 & 19 \\
                                   & Form & 4.16 & 3.82 & 1.0 & 3.0  & 6.0  & 1 & 15 \\
                                   & Ask  & 4.97 & 4.15 & 1.0 & 4.0  & 7.0  & 1 & 18 \\
                                   & Chat & 5.54 & 4.78 & 2.0 & 4.0  & 8.0  & 1 & 19 \\
\midrule
\multirow{4}{*}{Effort}            & Aggregated  & 6.13 & 4.57 & 3.0 & 5.0  & 10.0 & 1 & 20 \\
                                   & Form & 5.68 & 4.56 & 2.0 & 4.0  & 9.0  & 1 & 17 \\
                                   & Ask  & 6.39 & 4.41 & 3.0 & 5.0  & 10.0 & 1 & 19 \\
                                   & Chat & 6.32 & 4.77 & 3.0 & 5.0  & 10.0 & 1 & 20 \\
\midrule
\multirow{4}{*}{Frustration}       & Aggregated  & 6.86 & 5.24 & 2.5 & 5.0  & 11.0 & 1 & 21 \\
                                   & Form & 5.88 & 4.84 & 2.0 & 5.0  & 8.0  & 1 & 18 \\
                                   & Ask  & 7.58 & 5.34 & 3.0 & 6.0  & 12.0 & 1 & 21 \\
                                   & Chat & 7.13 & 5.44 & 3.0 & 5.0  & 11.0 & 1 & 19 \\
\midrule
\multirow{4}{*}{Performance (pos.)} & Aggregated  & 8.93 & 5.82 & 4.0 & 9.0  & 13.0 & 1 & 21 \\
                                    & Form & 8.71 & 5.48 & 4.0 & 8.0  & 12.0 & 1 & 21 \\
                                    & Ask  & 8.68 & 5.64 & 4.0 & 8.0  & 12.0 & 1 & 19 \\
                                    & Chat & 9.39 & 6.36 & 3.0 & 10.0 & 14.0 & 1 & 21 \\
\bottomrule
\end{tabular}
\caption{NASA R-TLX, full per-subscale descriptive distribution by
mode ($n = 69$ per mode), $1$--$21$ scale. Source:
\texttt{tlx\_per\_item\_summary.csv}.}
\label{tab:rtlx-per-subscale-distribution}
\end{table}

\begin{figure}[!htbp]
\centering
\includegraphics[width=0.45\linewidth]{rtlx_mental_demand.png}\hfill
\includegraphics[width=0.45\linewidth]{rtlx_physical_demand.png}\\[1ex]
\includegraphics[width=0.45\linewidth]{rtlx_temporal_demand.png}\hfill
\includegraphics[width=0.45\linewidth]{rtlx_effort.png}\\[1ex]
\includegraphics[width=0.45\linewidth]{rtlx_frustration.png}\hfill
\includegraphics[width=0.45\linewidth]{rtlx_performance.png}
\caption{Box plots of the six R-TLX subscales by mode
($n = 69$ per mode), one plot per subscale.}
\label{fig:rtlx-subquestions}
\end{figure}

\begin{figure}[!htbp]
\centering
\includegraphics[width=0.85\linewidth]{per_mode_timing_boxplot.png}
\caption{Per-mode completion time in seconds, extracted from the
\texttt{Time elapsed} line emitted by ROPAgen at document
submission ($n = 113$ documents: Form $37$, Ask $37$, Chat $39$).
Individual outliers are not drawn so the inter-quartile ranges are
legible.}
\label{fig:per-mode-timing}
\end{figure}

\subsection{Perceived AI Support}
\label{app:items-ai}

The compact per-item means table and the box plots are reproduced
inline with the Perceived AI Support results in
Section~\ref{sec:eval-ai-support}; the full per-item descriptive
distribution is reproduced here.

\begin{sidewaystable}
\centering
\small
\begin{tabularx}{\textheight}{@{}c X l r r r r r r r@{}}
\toprule
\# & Item & Mode & Mean & SD & Q1 & Median & Q3 & Min & Max \\
\midrule
\multirow{4}{*}{1} & \multirow{4}{*}{The AI was helpful in creating this document.}
                   & Aggregated  & 4.23 & 0.91 & 4 & 4 & 5 & 1 & 5 \\
                   & & Form & 4.23 & 0.97 & 4 & 5 & 5 & 1 & 5 \\
                   & & Ask  & 4.32 & 0.76 & 4 & 4 & 5 & 1 & 5 \\
                   & & Chat & 4.13 & 1.00 & 4 & 4 & 5 & 1 & 5 \\
\midrule
\multirow{4}{*}{2} & \multirow{4}{*}{The AI helped me to complete the document faster.}
                   & Aggregated  & 4.22 & 1.13 & 4 & 5 & 5 & 1 & 5 \\
                   & & Form & 4.43 & 0.96 & 4 & 5 & 5 & 1 & 5 \\
                   & & Ask  & 4.25 & 1.03 & 4 & 5 & 5 & 1 & 5 \\
                   & & Chat & 3.99 & 1.32 & 4 & 4 & 5 & 1 & 5 \\
\midrule
\multirow{4}{*}{3} & \multirow{4}{*}{I understood why the AI suggested specific GDPR fields.}
                   & Aggregated  & 3.81 & 1.06 & 3 & 4 & 5 & 1 & 5 \\
                   & & Form & 3.51 & 1.17 & 3 & 4 & 4 & 1 & 5 \\
                   & & Ask  & 4.16 & 0.92 & 4 & 4 & 5 & 1 & 5 \\
                   & & Chat & 3.75 & 0.98 & 3 & 4 & 5 & 1 & 5 \\
\midrule
\multirow{4}{*}{4} & \multirow{4}{*}{I trust that the AI made correct suggestions.}
                   & Aggregated  & 3.63 & 0.97 & 3 & 4 & 4 & 1 & 5 \\
                   & & Form & 3.41 & 1.09 & 3 & 4 & 4 & 1 & 5 \\
                   & & Ask  & 3.80 & 0.87 & 3 & 4 & 4 & 1 & 5 \\
                   & & Chat & 3.68 & 0.92 & 3 & 4 & 4 & 1 & 5 \\
\midrule
\multirow{4}{*}{5} & \multirow{4}{*}{I am confident that the AI identified the correct legal basis.}
                   & Aggregated  & 3.71 & 0.88 & 3 & 4 & 4 & 1 & 5 \\
                   & & Form & 3.58 & 0.90 & 3 & 4 & 4 & 1 & 5 \\
                   & & Ask  & 3.74 & 0.87 & 3 & 4 & 4 & 1 & 5 \\
                   & & Chat & 3.81 & 0.88 & 3 & 4 & 4 & 1 & 5 \\
\midrule
\multirow{4}{*}{6} & \multirow{4}{*}{I would use this mode again.}
                   & Aggregated  & 3.65 & 1.25 & 3 & 4 & 5 & 1 & 5 \\
                   & & Form & 3.99 & 1.14 & 4 & 4 & 5 & 1 & 5 \\
                   & & Ask  & 3.55 & 1.16 & 3 & 4 & 4 & 1 & 5 \\
                   & & Chat & 3.42 & 1.39 & 2 & 4 & 5 & 1 & 5 \\
\bottomrule
\end{tabularx}
\caption{Perceived AI Support, full per-item descriptive
distribution by mode ($n = 69$ per mode). Source:
\texttt{ai\_support\_per\_item\_summary.csv}.}
\label{tab:ai-per-item-distribution}
\end{sidewaystable}
\FloatBarrier

\begin{figure}[!htbp]
\centering
\includegraphics[width=0.45\linewidth]{ai_item_1_helpful.png}\hfill
\includegraphics[width=0.45\linewidth]{ai_item_2_efficiency.png}\\[1ex]
\includegraphics[width=0.45\linewidth]{ai_item_3_transparency.png}\hfill
\includegraphics[width=0.45\linewidth]{ai_item_4_trust.png}\\[1ex]
\includegraphics[width=0.45\linewidth]{ai_item_5_legal.png}\hfill
\includegraphics[width=0.45\linewidth]{ai_item_6_reuse.png}
\caption{Box plots of the six Perceived LLM Support items by mode
($n = 69$ per mode), one plot per item in administered order.}
\label{fig:ai-subquestions}
\end{figure}


\section{NLP Metric Scale Reference}
\label{app:metric-scales}

The five NLP metrics report on sharply different scales, so the
same numeral means very different things from one metric to the
next. They fall into two regimes. The lexical metrics (BLEU,
ROUGE, METEOR) span a wide range but sit low in absolute terms
when a candidate paraphrases the reference rather than reusing
its wording, with BLEU the most pessimistic of all because it
requires contiguous $n$-gram matches. The embedding-based
metrics (BERTScore, SBERT) are compressed toward the top of
the $[0,1]$ interval, because the contextual embeddings of any
two fluent texts on the same topic are already broadly similar;
for these, a difference of a few hundredths is substantial and
the absolute value should not be read on the same intuition as
a lexical score. Table~\ref{tab:metric-scales} gives heuristic
weak-versus-strong bands for the single-reference, free-text
comparison used in this thesis, alongside the mean actually
observed across the $113$ retained documents. The bands are
orientation aids for reading the results, not universal
thresholds: they are specific to comparison against one
expert-written reference, and the embedding bands in particular
are corpus-dependent.

\begin{table}[H]
\centering
\begin{tabularx}{\linewidth}{@{}l c c c c X@{}}
\toprule
Metric & Range & Weak ($\lesssim$) & Strong ($\gtrsim$) & Observed & Note \\
\midrule
BLEU              & $0$--$1$ & $0.05$ & $0.30$ & $0.13$ & Floor-heavy; needs verbatim $n$-grams \\
ROUGE-1           & $0$--$1$ & $0.30$ & $0.55$ & $0.49$ & Unigram recall \\
ROUGE-2           & $0$--$1$ & $0.15$ & $0.35$ & $0.29$ & Bigrams; structurally lower \\
ROUGE-L           & $0$--$1$ & $0.25$ & $0.45$ & $0.38$ & Longest common subsequence \\
METEOR            & $0$--$1$ & $0.20$ & $0.50$ & $0.39$ & Synonym and stem credit \\
\midrule
BERTScore P/R/F1  & $0$--$1$ & $0.82$ & $0.92$ & $0.90$ & Compressed; below $0.80$ rare \\
SBERT cosine      & $0$--$1$ & $0.90$ & $0.97$ & $0.99$ & Most compressed; above $0.95$ is baseline \\
\bottomrule
\end{tabularx}
\caption{Heuristic scale reference for the NLP metrics. Weak and
strong columns are approximate orientation bands for free-text
comparison against a single reference, not universal thresholds;
the observed column is the aggregated mean across the $113$
retained documents (Table~\ref{tab:nlp-doc-overview}). Lexical
metrics span a wide range but read low for paraphrased content;
embedding-based metrics are compressed near the top of their
range, so small differences are meaningful.}
\label{tab:metric-scales}
\end{table}

\section{Worked Examples (Document Quality)}
\label{app:worked-examples}

These passage-level worked examples were referenced from earlier
drafts of the Document Quality chapter and from the Discussion
case-study framing, and are preserved here in case they are
returned to. They contrast Participant A (Chat) and Participant
B (Form) passages on Retention (Chapter~5) and Purpose
(Chapter~2) with per-passage NLP metric values.

\paragraph{Document-scale reading.} Two participants make the
confidence-quality gap visible at document scale. One produced
the study's strongest document on every measure; the other
produced a complete-looking document filled with details the
scenario never mentioned, which the NLP pipeline did not catch.
Participant A's Chat document is the case-study ceiling on both
NLP and judge readings: highest semantic similarity to the
reference, highest scenario faithfulness on the judge rubric,
and the participant's strongest output across the metric suite.
Participant B's Form document is the case-study floor on the
judge rubric: maximum Completeness, lowest Faithfulness and
Hallucination-inverse, with invented financial data, biometric
data, video surveillance, and named technical standards absent
from the scenario. The Retention and Purpose passages reproduced
in Figures~\ref{fig:worked-lexical-a}--\ref{fig:worked-semantic-b}
show both participants at passage scale.

\paragraph{Passage-scale reading.} At passage scale the same
pattern shows up. On Retention and Purpose alike, Participant B's
Form passages outscored Participant A's Chat passages on every
reported metric, even though Chat's content was faithful and
Form's was fabricated. The Retention passage cited $6$ to $10$
year statutory retention against the scenario's $3$-month rule;
the Purpose passage cited legal articles (\S~257~HGB,
\S~147~AO, \S~239~Abs.~1~HGB, \S~2~Abs.~1~Nr.~1~BDSG) absent
from the scenario. BLEU illustrates the limit sharpest:
Participant A's Chat Retention passage shares zero $4$-grams
with the reference, so BLEU-4 zeros at the $4$-gram tier
independent of the brevity penalty.

\subsection{Lexical Metrics: Retention Chapter}
\label{app:worked-lexical}

To attach the lexical metrics to concrete text,
Figures~\ref{fig:worked-lexical-a} and~\ref{fig:worked-lexical-b}
show the Retention chapter (Chapter~5,
\emph{Aufbewahrungsfristen und Löschung}) for two case-study
participants alongside the reference passage and the
corresponding per-passage metric values. Chapter~5 is used here
because it has the lowest absolute lexical scores and the widest
variance of any chapter, and the case-study Form passage there
carries the most legible fabrications, which makes the lexical
metrics' blindness to truth easiest to see.

\begingroup
\lstset{
  basicstyle=\ttfamily\scriptsize,
  breaklines=true,
  breakatwhitespace=false,
  breakindent=0pt,
  postbreak=\mbox{$\hookrightarrow$\space},
  frame=single,
  framerule=0.4pt,
  language=,
  numbers=none,
  showstringspaces=false,
  keepspaces=true,
  columns=fullflexible,
  upquote=true,
  aboveskip=0.4em,
  belowskip=0.4em,
  literate=%
    {ä}{{\"a}}1 {ö}{{\"o}}1 {ü}{{\"u}}1
    {Ä}{{\"A}}1 {Ö}{{\"O}}1 {Ü}{{\"U}}1 {ß}{{\ss}}1
    {§}{{\S}}1
    {•}{{\textbullet}}1
    {–}{{--}}1 {—}{{---}}1
    {“}{{``}}1 {”}{{''}}1
    {‘}{{`}}1 {’}{{'}}1
}
\begin{figure}[!htbp]
\centering
\noindent\textit{Reference passage (single-sentence statement of the scenario's 3-month rule).}
\begin{lstlisting}
Wir löschen die personenbezogenen Daten drei Monate nach Beendigung des Beschäftigungsverhältnisses, sofern keine gesetzlichen Aufbewahrungspflichten entgegenstehen.
\end{lstlisting}

\noindent\textit{Participant A (Chat) passage.}
\begin{lstlisting}
### 5.1 Löschfristen
Die personenbezogenen Daten werden gemäß folgender Regelungen gelöscht:

- **Grundsatz:** Die Daten werden spätestens **3 Monate nach Austritt der Mitarbeitenden** gelöscht.
- **Ausnahmen:** Sofern gesetzliche Aufbewahrungspflichten (z. B. im Rahmen laufender Verfahren) bestehen, erfolgt die Löschung erst nach Ablauf dieser Fristen.

### 5.2 Verfahren zur Löschung
Die Löschung erfolgt entweder automatisiert durch unsere Systeme oder manuell durch die zuständigen Mitarbeitenden. Die Einhaltung der Löschfristen wird regelmäßig überprüft.
\end{lstlisting}
\caption{Worked example for the lexical metrics (1/2): reference
and Participant A (Chat) Retention passages (Chapter~5,
\emph{Aufbewahrungsfristen und L\"oschung}). Continues in
Figure~\ref{fig:worked-lexical-b}.}
\label{fig:worked-lexical-a}
\end{figure}

\begin{figure}[!htbp]
\centering
\noindent\textit{Participant B (Form) passage.}
\begin{lstlisting}
Wir bewahren personenbezogene Daten unserer Mitarbeitenden gemäß den gesetzlichen Vorgaben auf und löschen diese nach Ablauf der festgelegten Fristen. Die Aufbewahrungsfristen richten sich nach den folgenden Grundsätzen:

- **Grundsatz:** 6 Jahre nach Beendigung des Beschäftigungsverhältnisses (gemäß § 257 HGB und § 147 AO für steuerrelevante Unterlagen).
- **Ausnahmen:**
  - Sozialversicherungsnummern und Vertragsdaten können bis zu 10 Jahre aufbewahrt werden, sofern dies aufgrund rechtlicher Auseinandersetzungen oder Verdachtsmomenten (z. B. Sozialversicherungsbetrug) erforderlich ist.
  - Biometrische Daten und Videoüberwachungsdaten werden nur so lange gespeichert, wie dies für den jeweiligen Zweck (z. B. Zutrittskontrolle) notwendig ist.

Nach Ablauf der Aufbewahrungsfristen werden die Daten automatisch gelöscht, sofern keine gesetzlichen oder vertraglichen Gründe für eine längere Speicherung vorliegen.
\end{lstlisting}

\vspace{0.5ex}
\begin{tabular}{@{}l c c c@{}}
\toprule
Candidate    & BLEU   & ROUGE-2 & METEOR \\
\midrule
Participant A (Chat) & 0.0000 & 0.0625  & 0.1724 \\
Participant B (Form) & \textbf{0.0356} & \textbf{0.1185}  & \textbf{0.3790} \\
\bottomrule
\end{tabular}
\caption{Worked example for the lexical metrics (2/2):
Participant B (Form) Retention passage with per-passage BLEU, ROUGE-2,
and METEOR values. Per-passage metric values are computed with
the same pipeline as the chapter-level table; see
Appendix~\ref{app:items-chapter-nlp-lexical} for the full chapter
means. Continued from Figure~\ref{fig:worked-lexical-a}.}
\label{fig:worked-lexical-b}
\end{figure}
\endgroup

On the Retention passages above, Participant A's Chat passage
scores BLEU $= 0.0000$ against the reference's single-sentence
statement (the value the metric assigns when no candidate
$4$-gram appears in the reference) and Participant B's Form
passage scores BLEU $= 0.0356$, where a handful of shared
$4$-grams survive. The ROUGE-2 and METEOR readings move in the
same direction: Participant A's Chat passage scores ROUGE-2
$= 0.0625$ and METEOR $= 0.1724$, while Participant B's Form
passage scores ROUGE-2 $= 0.1185$ on the strength of bigrams
shared with the reference and METEOR $= 0.3790$ once stems and
WordNet synonyms are credited.

\subsection{Embedding-Based Metrics: Purpose Chapter}
\label{app:worked-semantic}

To attach the embedding-based metrics to concrete text,
Figures~\ref{fig:worked-semantic-a} and~\ref{fig:worked-semantic-b}
show the Purpose chapter (Chapter~2, \emph{Zwecke und
Rechtsgrundlagen der Verarbeitung}) for the same two case-study
participants alongside the reference passage and the
corresponding per-passage metric values. Chapter~2 is used here
because both case-study passages cover the same purposes
faithfully but in different wording, which is the case the
embedding-based metrics are designed to credit.

\begingroup
\lstset{
  basicstyle=\ttfamily\scriptsize,
  breaklines=true,
  breakatwhitespace=false,
  breakindent=0pt,
  postbreak=\mbox{$\hookrightarrow$\space},
  frame=single,
  framerule=0.4pt,
  language=,
  numbers=none,
  showstringspaces=false,
  keepspaces=true,
  columns=fullflexible,
  upquote=true,
  aboveskip=0.4em,
  belowskip=0.4em,
  literate=%
    {ä}{{\"a}}1 {ö}{{\"o}}1 {ü}{{\"u}}1
    {Ä}{{\"A}}1 {Ö}{{\"O}}1 {Ü}{{\"U}}1 {ß}{{\ss}}1
    {§}{{\S}}1
    {•}{{\textbullet}}1
    {–}{{--}}1 {—}{{---}}1
    {“}{{``}}1 {”}{{''}}1
    {‘}{{`}}1 {’}{{'}}1
}
\begin{figure}[!htbp]
\centering
\noindent\textit{Reference passage (purposes plus the two legal bases for processing).}
\begin{lstlisting}
2.1 Verarbeitungszwecke
Wir verarbeiten personenbezogene Daten für folgende Zwecke:
• Verwaltung von Mitarbeitenden vom Eintritt bis zum Austritt, einschließlich:
• Anlegen neuer Mitarbeitender
• Dokumentation von Arbeitszeiten
• Verwaltung von Urlaub und Abwesenheiten
• Verwaltung von Parkplatzberechtigungen für berechtigte Mitarbeitende

2.2 Rechtsgrundlagen der Verarbeitung
Die Verarbeitung erfolgt auf Grundlage folgender Rechtsnormen:
• Art. 6 Abs. 1 lit. b DSGVO (Erfüllung eines Vertrags oder vorvertraglicher Maßnahmen)
• § 26 BDSG (Datenverarbeitung für Zwecke des Beschäftigungsverhältnisses)
\end{lstlisting}

\noindent\textit{Participant A (Chat) passage.}
\begin{lstlisting}
### 2.1 Zwecke der Verarbeitung
Wir verarbeiten personenbezogene Daten unserer Mitarbeitenden zur ganzheitlichen Verwaltung von deren Beschäftigungsverhältnis – vom Eintritt bis zum Austritt aus dem Unternehmen. Die Verarbeitung dient folgenden Zwecken:

- Erfassung und Pflege von Stammdaten neuer und bestehender Mitarbeitender,
- Dokumentation von Arbeitszeiten sowie Verwaltung von Urlaubs- und Abwesenheitszeiten,
- Zuweisung und Verwaltung von Parkplatzberechtigungen für den firmeneigenen Parkplatz mit Schranke. Hierfür wird ein eindeutiges Merkmal (Fahrzeugkennzeichen) hinterlegt, um die Nutzung zu ermöglichen.

Die Verarbeitung erfolgt ausschließlich auf internen Systemen (Microsoft Office 365 und eigene Server) ohne Nutzung externer Cloud-Dienste. Medizinische Daten oder Krankentage sind nicht Gegenstand dieser Verarbeitungstätigkeit.

### 2.2 Rechtsgrundlagen der Verarbeitung
Die Verarbeitung personenbezogener Daten erfolgt auf folgenden Rechtsgrundlagen gemäß Art. 6 Abs. 1 DSGVO sowie ergänzenden nationalen Vorschriften:

- **Art. 6 Abs. 1 lit. b DSGVO (Vertragserfüllung):**
  Die Verarbeitung ist für die Durchführung des Beschäftigungsverhältnisses erforderlich, insbesondere zur Verwaltung von Arbeitszeiten, Urlaub und Parkplatzberechtigungen.

- **Art. 6 Abs. 1 lit. f DSGVO (Berechtigtes Interesse):**
  Die Verarbeitung dient unserem berechtigten Interesse an einer effizienten Verwaltung der Mitarbeiterressourcen sowie der Sicherstellung einer geordneten Nutzung des firmeneigenen Parkplatzes.

- **§ 26 BDSG (Datenverarbeitung für Zwecke des Beschäftigungsverhältnisses):**
  Die Verarbeitung erfolgt im Rahmen des Beschäftigungsverhältnisses und ist für dessen Durchführung erforderlich.
\end{lstlisting}

\caption{Worked example for the embedding-based metrics (1/2):
reference and Participant A (Chat) Purpose passages (Chapter~2,
\emph{Zwecke und Rechtsgrundlagen der Verarbeitung}). Continues
in Figure~\ref{fig:worked-semantic-b}.}
\label{fig:worked-semantic-a}
\end{figure}

\begin{figure}[!htbp]
\centering
\noindent\textit{Participant B (Form) passage.}
\begin{lstlisting}
Wir verarbeiten personenbezogene Daten unserer Mitarbeitenden ausschließlich zur Verwaltung des Beschäftigungsverhältnisses von der Einstellung bis zum Austritt. Dies umfasst folgende Zwecke:

- Erfassung und Pflege von Stammdaten der Mitarbeitenden
- Dokumentation von Arbeitszeiten und Anwesenheiten
- Verwaltung von Urlaubs- und Abwesenheitszeiten
- Verwaltung von Parkplatzberechtigungen (z. B. Fahrzeugkennzeichen)
- Erfüllung gesetzlicher Aufbewahrungs- und Meldepflichten
- Sicherstellung der betrieblichen Abläufe und Compliance

Die Verarbeitung erfolgt auf Grundlage folgender Rechtsgrundlagen:

- **Art. 6 Abs. 1 lit. b DSGVO** (Erfüllung eines Vertrags oder vorvertraglicher Maßnahmen)
- **§ 26 BDSG** (Datenverarbeitung für Zwecke des Beschäftigungsverhältnisses)
- **§ 257 HGB** (Aufbewahrung von Handelsbüchern und Unterlagen)
- **§ 147 AO** (Aufbewahrung von steuerrelevanten Unterlagen)
- **§ 239 Abs. 1 HGB** (Führung der Handelsbücher)
- **§ 2 Abs. 1 Nr. 1 BDSG** (Verarbeitung personenbezogener Daten im Beschäftigungskontext)
\end{lstlisting}

\vspace{0.5ex}
\begin{tabular}{@{}l c c c c@{}}
\toprule
Candidate    & BERTScore P & BERTScore R & BERTScore F1 & SBERT \\
\midrule
Participant A (Chat) & 0.8797      & 0.9180      & 0.8984       & 0.9810 \\
Participant B (Form) & \textbf{0.8954}      & \textbf{0.9336}      & \textbf{0.9141}       & \textbf{0.9882} \\
\bottomrule
\end{tabular}
\caption{Worked example for the embedding-based metrics (2/2):
Participant B (Form) Purpose passage with per-passage BERTScore P/R/F1
and SBERT values. SBERT cosine similarity is reported under the
ModernBERT encoder, matching the document- and chapter-level
tables. Per-passage metric values are computed with the same
pipeline as the chapter-level table; see
Appendix~\ref{app:items-chapter-nlp-bert} for the full chapter
means. Continued from Figure~\ref{fig:worked-semantic-a}.}
\label{fig:worked-semantic-b}
\end{figure}
\endgroup

On the Purpose passages above, Participant A's Chat passage
scores BERTScore Precision $= 0.8797$, Recall $= 0.9180$, and F1
$= 0.8984$ against the reference, while Participant B's Form
passage scores Precision $= 0.8954$, Recall $= 0.9336$, and F1
$= 0.9141$, with the embedding model crediting its broader
overlap with the reference's legal vocabulary. SBERT is the more
forgiving metric of the two: Participant A's Chat passage scores
SBERT $= 0.9810$ and Participant B's Form passage scores SBERT
$= 0.9882$, with the document-level encoder placing both
passages near the top of the metric's range.


\section{Document-Level NLP Metric Distributions}
\label{app:items-nlp-doc}

The compact summaries and the box plots are reproduced inline with
the NLP results in Section~\ref{sec:eval-nlp-doc}; the full
per-document descriptive distributions are reproduced here, split
into the lexical and embedding-based families.

\subsection{Lexical Metrics}
\label{app:items-nlp-doc-lexical}

Table~\ref{tab:nlp-doc-distrib-lexical} reports the per-mode mean,
dispersion, and quartiles for each surface-overlap metric on the
full retained set of $113$ documents.

\begin{sidewaystable}
\centering
\small
\begin{tabular}{@{}l l r r r r r r r r@{}}
\toprule
Metric & Mode & $n$ & Mean & SD & Q1 & Median & Q3 & Min & Max \\
\midrule
\multirow{4}{*}{BLEU}    & Aggregated  & 113 & 0.133 & 0.048 & 0.094 & 0.128 & 0.159 & 0.047 & 0.246 \\
                         & Form & 37  & 0.129 & 0.048 & 0.088 & 0.124 & 0.158 & 0.064 & 0.230 \\
                         & Ask  & 37  & 0.127 & 0.056 & 0.083 & 0.125 & 0.148 & 0.047 & 0.246 \\
                         & Chat & 39  & 0.143 & 0.039 & 0.122 & 0.134 & 0.163 & 0.067 & 0.235 \\
\midrule
\multirow{4}{*}{ROUGE-1} & Aggregated  & 113 & 0.488 & 0.087 & 0.425 & 0.491 & 0.560 & 0.313 & 0.645 \\
                         & Form & 37  & 0.488 & 0.098 & 0.421 & 0.493 & 0.569 & 0.313 & 0.644 \\
                         & Ask  & 37  & 0.465 & 0.088 & 0.412 & 0.449 & 0.522 & 0.322 & 0.618 \\
                         & Chat & 39  & 0.511 & 0.071 & 0.468 & 0.517 & 0.561 & 0.330 & 0.645 \\
\midrule
\multirow{4}{*}{ROUGE-2} & Aggregated  & 113 & 0.285 & 0.077 & 0.234 & 0.268 & 0.344 & 0.121 & 0.455 \\
                         & Form & 37  & 0.306 & 0.076 & 0.241 & 0.301 & 0.361 & 0.178 & 0.455 \\
                         & Ask  & 37  & 0.257 & 0.080 & 0.200 & 0.238 & 0.275 & 0.121 & 0.429 \\
                         & Chat & 39  & 0.291 & 0.069 & 0.253 & 0.283 & 0.344 & 0.131 & 0.447 \\
\midrule
\multirow{4}{*}{ROUGE-L} & Aggregated  & 113 & 0.375 & 0.087 & 0.311 & 0.359 & 0.441 & 0.194 & 0.566 \\
                         & Form & 37  & 0.388 & 0.094 & 0.320 & 0.391 & 0.468 & 0.238 & 0.566 \\
                         & Ask  & 37  & 0.347 & 0.091 & 0.296 & 0.326 & 0.376 & 0.194 & 0.545 \\
                         & Chat & 39  & 0.389 & 0.072 & 0.348 & 0.387 & 0.444 & 0.230 & 0.544 \\
\midrule
\multirow{4}{*}{METEOR}  & Aggregated  & 113 & 0.389 & 0.062 & 0.359 & 0.388 & 0.431 & 0.234 & 0.514 \\
                         & Form & 37  & 0.412 & 0.051 & 0.367 & 0.419 & 0.442 & 0.335 & 0.514 \\
                         & Ask  & 37  & 0.365 & 0.068 & 0.327 & 0.366 & 0.409 & 0.234 & 0.488 \\
                         & Chat & 39  & 0.390 & 0.057 & 0.374 & 0.389 & 0.421 & 0.235 & 0.501 \\
\bottomrule
\end{tabular}
\caption{Document-level lexical metric distributions by mode (full
sample $n = 113$). Source: \texttt{doc\_nlp\_summary.csv}.}
\label{tab:nlp-doc-distrib-lexical}
\end{sidewaystable}
\FloatBarrier

\begin{figure}[!htbp]
\centering
\includegraphics[width=0.45\linewidth]{nlp_doc_bleu.png}\hfill
\includegraphics[width=0.45\linewidth]{nlp_doc_rouge1.png}\\[1ex]
\includegraphics[width=0.45\linewidth]{nlp_doc_rouge2.png}\hfill
\includegraphics[width=0.45\linewidth]{nlp_doc_rougeL.png}\\[1ex]
\includegraphics[width=0.45\linewidth]{nlp_doc_meteor.png}
\caption{Document-level distributions of BLEU, ROUGE-1, ROUGE-2,
ROUGE-L, and METEOR by mode (full sample $n = 113$).}
\label{fig:nlp-doc-boxplot}
\end{figure}

\subsection{BERT-Based and Sentence-Embedding Metrics}
\label{app:items-nlp-doc-bert}

Table~\ref{tab:nlp-doc-distrib-bert} reports the same descriptive
breakdown for the embedding-based metrics across the three modes.

\begin{sidewaystable}
\centering
\small
\begin{tabular}{@{}l l r r r r r r r r@{}}
\toprule
Metric & Mode & $n$ & Mean & SD & Q1 & Median & Q3 & Min & Max \\
\midrule
\multirow{4}{*}{BERTScore Precision} & Aggregated  & 113 & 0.890 & 0.014 & 0.880 & 0.891 & 0.900 & 0.861 & 0.922 \\
                                     & Form & 37  & 0.889 & 0.014 & 0.880 & 0.889 & 0.899 & 0.862 & 0.912 \\
                                     & Ask  & 37  & 0.886 & 0.016 & 0.872 & 0.889 & 0.899 & 0.861 & 0.918 \\
                                     & Chat & 39  & 0.895 & 0.012 & 0.890 & 0.896 & 0.902 & 0.867 & 0.922 \\
\midrule
\multirow{4}{*}{BERTScore Recall}    & Aggregated  & 113 & 0.915 & 0.014 & 0.914 & 0.918 & 0.922 & 0.861 & 0.936 \\
                                     & Form & 37  & 0.923 & 0.005 & 0.919 & 0.923 & 0.926 & 0.914 & 0.936 \\
                                     & Ask  & 37  & 0.908 & 0.017 & 0.908 & 0.915 & 0.918 & 0.861 & 0.925 \\
                                     & Chat & 39  & 0.913 & 0.013 & 0.908 & 0.917 & 0.922 & 0.874 & 0.933 \\
\midrule
\multirow{4}{*}{BERTScore F1}        & Aggregated  & 113 & 0.902 & 0.011 & 0.895 & 0.904 & 0.910 & 0.863 & 0.922 \\
                                     & Form & 37  & 0.905 & 0.009 & 0.900 & 0.906 & 0.912 & 0.887 & 0.920 \\
                                     & Ask  & 37  & 0.897 & 0.013 & 0.892 & 0.896 & 0.906 & 0.863 & 0.918 \\
                                     & Chat & 39  & 0.904 & 0.010 & 0.899 & 0.906 & 0.911 & 0.872 & 0.922 \\
\midrule
\multirow{4}{*}{SBERT (ModernBERT)}  & Aggregated  & 113 & 0.987 & 0.004 & 0.985 & 0.988 & 0.990 & 0.971 & 0.994 \\
                                     & Form & 37  & 0.986 & 0.004 & 0.984 & 0.988 & 0.989 & 0.974 & 0.992 \\
                                     & Ask  & 37  & 0.986 & 0.005 & 0.983 & 0.987 & 0.990 & 0.971 & 0.993 \\
                                     & Chat & 39  & 0.990 & 0.002 & 0.988 & 0.990 & 0.991 & 0.985 & 0.994 \\
\bottomrule
\end{tabular}
\caption{Document-level BERT-based and SBERT distributions by mode
(full sample $n = 113$). Source: \texttt{doc\_bert\_summary.csv}.}
\label{tab:nlp-doc-distrib-bert}
\end{sidewaystable}
\FloatBarrier

\begin{figure}[!htbp]
\centering
\includegraphics[width=0.45\linewidth]{nlp_doc_bertscore_p.png}\hfill
\includegraphics[width=0.45\linewidth]{nlp_doc_bertscore_r.png}\\[1ex]
\includegraphics[width=0.45\linewidth]{nlp_doc_bertscore_f1.png}\hfill
\includegraphics[width=0.45\linewidth]{nlp_doc_sbert.png}
\caption{Document-level distributions of BERTScore Precision,
Recall, F1, and SBERT cosine similarity (ModernBERT) by mode
(full sample $n = 113$), one plot per metric.}
\label{fig:bert-doc-boxplot}
\end{figure}


\section{Chapter-Level NLP Metric Distributions}
\label{app:items-chapter-nlp}

The chapter-level box plots are reproduced inline with the
chapter-level NLP results in Section~\ref{sec:eval-nlp-chapter};
the per-chapter, per-mode means (with SDs in parentheses) for
every metric are reproduced here.

\subsection{Lexical Metrics}
\label{app:items-chapter-nlp-lexical}

\begin{table}[H]
\centering
\small
\setlength{\tabcolsep}{4pt}
\begin{tabular}{@{}l c c c c@{}}
\toprule
Chapter & Aggregated & Form & Ask & Chat \\
\midrule
Controller \& Contact         & 0.181 (0.080) & 0.179 (0.077) & 0.200 (0.085) & 0.166 (0.076) \\
Purposes \& Legal Bases       & 0.121 (0.060) & 0.145 (0.057) & 0.098 (0.038) & 0.121 (0.070) \\
Personal Data Categories      & 0.081 (0.050) & 0.110 (0.041) & 0.077 (0.052) & 0.058 (0.043) \\
Data Subjects \& Recipients   & 0.097 (0.075) & 0.123 (0.076) & 0.081 (0.082) & 0.088 (0.063) \\
Retention \& Deletion         & 0.044 (0.076) & 0.038 (0.044) & 0.023 (0.073) & 0.068 (0.094) \\
Technical \& Org. Measures    & 0.080 (0.066) & 0.078 (0.068) & 0.070 (0.080) & 0.091 (0.047) \\
\bottomrule
\end{tabular}
\caption{Chapter-level BLEU mean (SD) by mode (full sample
$n = 113$).}
\label{tab:nlp-chapter-bleu}
\end{table}

\begin{table}[H]
\centering
\small
\setlength{\tabcolsep}{4pt}
\begin{tabular}{@{}l c c c c@{}}
\toprule
Chapter & Aggregated & Form & Ask & Chat \\
\midrule
Controller \& Contact         & 0.571 (0.114) & 0.550 (0.097) & 0.605 (0.129) & 0.559 (0.109) \\
Purposes \& Legal Bases       & 0.414 (0.109) & 0.457 (0.094) & 0.373 (0.093) & 0.410 (0.122) \\
Personal Data Categories      & 0.467 (0.164) & 0.600 (0.084) & 0.366 (0.120) & 0.437 (0.174) \\
Data Subjects \& Recipients   & 0.437 (0.116) & 0.491 (0.111) & 0.393 (0.128) & 0.426 (0.086) \\
Retention \& Deletion         & 0.260 (0.131) & 0.252 (0.092) & 0.212 (0.132) & 0.314 (0.143) \\
Technical \& Org. Measures    & 0.331 (0.130) & 0.313 (0.137) & 0.330 (0.157) & 0.350 (0.090) \\
\bottomrule
\end{tabular}
\caption{Chapter-level ROUGE-1 mean (SD) by mode (full sample
$n = 113$).}
\label{tab:nlp-chapter-rouge1}
\end{table}

\begin{table}[H]
\centering
\small
\setlength{\tabcolsep}{4pt}
\begin{tabular}{@{}l c c c c@{}}
\toprule
Chapter & Aggregated & Form & Ask & Chat \\
\midrule
Controller \& Contact         & 0.386 (0.145) & 0.363 (0.131) & 0.430 (0.158) & 0.365 (0.138) \\
Purposes \& Legal Bases       & 0.269 (0.109) & 0.332 (0.093) & 0.216 (0.077) & 0.261 (0.119) \\
Personal Data Categories      & 0.298 (0.148) & 0.425 (0.097) & 0.198 (0.089) & 0.272 (0.149) \\
Data Subjects \& Recipients   & 0.250 (0.109) & 0.317 (0.105) & 0.215 (0.113) & 0.220 (0.077) \\
Retention \& Deletion         & 0.131 (0.113) & 0.134 (0.071) & 0.087 (0.107) & 0.168 (0.137) \\
Technical \& Org. Measures    & 0.176 (0.118) & 0.166 (0.120) & 0.163 (0.141) & 0.197 (0.089) \\
\bottomrule
\end{tabular}
\caption{Chapter-level ROUGE-2 mean (SD) by mode (full sample
$n = 113$).}
\label{tab:nlp-chapter-rouge2}
\end{table}

\begin{table}[H]
\centering
\small
\setlength{\tabcolsep}{4pt}
\begin{tabular}{@{}l c c c c@{}}
\toprule
Chapter & Aggregated & Form & Ask & Chat \\
\midrule
Controller \& Contact         & 0.534 (0.129) & 0.518 (0.110) & 0.576 (0.140) & 0.511 (0.128) \\
Purposes \& Legal Bases       & 0.351 (0.111) & 0.404 (0.096) & 0.299 (0.083) & 0.351 (0.125) \\
Personal Data Categories      & 0.426 (0.168) & 0.567 (0.085) & 0.309 (0.107) & 0.403 (0.179) \\
Data Subjects \& Recipients   & 0.372 (0.121) & 0.433 (0.114) & 0.330 (0.131) & 0.355 (0.096) \\
Retention \& Deletion         & 0.229 (0.136) & 0.218 (0.087) & 0.183 (0.141) & 0.283 (0.152) \\
Technical \& Org. Measures    & 0.257 (0.135) & 0.237 (0.140) & 0.246 (0.158) & 0.285 (0.100) \\
\bottomrule
\end{tabular}
\caption{Chapter-level ROUGE-L mean (SD) by mode (full sample
$n = 113$).}
\label{tab:nlp-chapter-rougeL}
\end{table}

\begin{table}[H]
\centering
\small
\setlength{\tabcolsep}{4pt}
\begin{tabular}{@{}l c c c c@{}}
\toprule
Chapter & Aggregated & Form & Ask & Chat \\
\midrule
Controller \& Contact         & 0.471 (0.128) & 0.456 (0.120) & 0.506 (0.136) & 0.453 (0.124) \\
Purposes \& Legal Bases       & 0.400 (0.113) & 0.452 (0.075) & 0.348 (0.101) & 0.401 (0.131) \\
Personal Data Categories      & 0.315 (0.120) & 0.415 (0.074) & 0.275 (0.111) & 0.258 (0.104) \\
Data Subjects \& Recipients   & 0.355 (0.109) & 0.440 (0.098) & 0.312 (0.101) & 0.314 (0.077) \\
Retention \& Deletion         & 0.314 (0.160) & 0.332 (0.114) & 0.245 (0.160) & 0.362 (0.177) \\
Technical \& Org. Measures    & 0.293 (0.113) & 0.295 (0.108) & 0.280 (0.145) & 0.305 (0.079) \\
\bottomrule
\end{tabular}
\caption{Chapter-level METEOR mean (SD) by mode (full sample
$n = 113$).}
\label{tab:nlp-chapter-meteor}
\end{table}

\begin{figure}[!htbp]
\centering
\includegraphics[width=0.7\linewidth]{chapter_nlp_rouge1.png}\\[1ex]
\includegraphics[width=0.7\linewidth]{chapter_nlp_rouge2.png}\\[1ex]
\includegraphics[width=0.7\linewidth]{chapter_nlp_rougeL.png}
\caption{Chapter-level distributions of ROUGE-1, ROUGE-2, and
ROUGE-L by mode (full sample $n = 113$), one plot per metric
showing all six ROPA chapters side by side.}
\label{fig:chapter-nlp-boxplot-a}
\end{figure}

\begin{figure}[!htbp]
\centering
\includegraphics[width=0.85\linewidth]{chapter_nlp_bleu.png}\\[1ex]
\includegraphics[width=0.85\linewidth]{chapter_nlp_meteor.png}
\caption{Chapter-level distributions of BLEU and METEOR by
mode (full sample $n = 113$), one plot per metric showing all
six ROPA chapters side by side.}
\label{fig:chapter-nlp-boxplot-b}
\end{figure}

\subsection{BERT-Based and Sentence-Embedding Metrics}
\label{app:items-chapter-nlp-bert}

\begin{table}[H]
\centering
\small
\setlength{\tabcolsep}{4pt}
\begin{tabular}{@{}l c c c c@{}}
\toprule
Chapter & Aggregated & Form & Ask & Chat \\
\midrule
Controller \& Contact         & 0.857 (0.017) & 0.855 (0.015) & 0.862 (0.020) & 0.855 (0.016) \\
Purposes \& Legal Bases       & 0.873 (0.021) & 0.872 (0.020) & 0.869 (0.020) & 0.877 (0.022) \\
Personal Data Categories      & 0.883 (0.022) & 0.897 (0.015) & 0.868 (0.019) & 0.884 (0.022) \\
Data Subjects \& Recipients   & 0.857 (0.028) & 0.856 (0.025) & 0.853 (0.035) & 0.864 (0.023) \\
Retention \& Deletion         & 0.824 (0.042) & 0.814 (0.033) & 0.817 (0.050) & 0.840 (0.037) \\
Technical \& Org. Measures    & 0.854 (0.033) & 0.841 (0.038) & 0.858 (0.031) & 0.863 (0.027) \\
\bottomrule
\end{tabular}
\caption{Chapter-level BERTScore Precision mean (SD) by mode (full
sample $n = 113$).}
\label{tab:nlp-chapter-bertp}
\end{table}

\begin{table}[H]
\centering
\small
\setlength{\tabcolsep}{4pt}
\begin{tabular}{@{}l c c c c@{}}
\toprule
Chapter & Aggregated & Form & Ask & Chat \\
\midrule
Controller \& Contact         & 0.882 (0.022) & 0.879 (0.019) & 0.887 (0.025) & 0.879 (0.022) \\
Purposes \& Legal Bases       & 0.902 (0.024) & 0.911 (0.012) & 0.892 (0.028) & 0.904 (0.024) \\
Personal Data Categories      & 0.878 (0.041) & 0.910 (0.016) & 0.862 (0.042) & 0.861 (0.039) \\
Data Subjects \& Recipients   & 0.880 (0.026) & 0.904 (0.014) & 0.873 (0.024) & 0.864 (0.022) \\
Retention \& Deletion         & 0.889 (0.024) & 0.890 (0.018) & 0.879 (0.023) & 0.897 (0.027) \\
Technical \& Org. Measures    & 0.888 (0.025) & 0.889 (0.018) & 0.882 (0.032) & 0.893 (0.022) \\
\bottomrule
\end{tabular}
\caption{Chapter-level BERTScore Recall mean (SD) by mode (full
sample $n = 113$).}
\label{tab:nlp-chapter-bertr}
\end{table}

\begin{table}[H]
\centering
\small
\setlength{\tabcolsep}{4pt}
\begin{tabular}{@{}l c c c c@{}}
\toprule
Chapter & Aggregated & Form & Ask & Chat \\
\midrule
Controller \& Contact         & 0.869 (0.019) & 0.867 (0.016) & 0.875 (0.022) & 0.867 (0.017) \\
Purposes \& Legal Bases       & 0.887 (0.018) & 0.891 (0.015) & 0.880 (0.018) & 0.890 (0.020) \\
Personal Data Categories      & 0.880 (0.029) & 0.903 (0.012) & 0.865 (0.025) & 0.872 (0.029) \\
Data Subjects \& Recipients   & 0.868 (0.020) & 0.879 (0.017) & 0.862 (0.022) & 0.863 (0.017) \\
Retention \& Deletion         & 0.855 (0.031) & 0.850 (0.024) & 0.847 (0.035) & 0.867 (0.031) \\
Technical \& Org. Measures    & 0.871 (0.025) & 0.864 (0.027) & 0.870 (0.027) & 0.878 (0.019) \\
\bottomrule
\end{tabular}
\caption{Chapter-level BERTScore F1 mean (SD) by mode (full sample
$n = 113$).}
\label{tab:nlp-chapter-bertf1}
\end{table}

\begin{table}[H]
\centering
\small
\setlength{\tabcolsep}{4pt}
\begin{tabular}{@{}l c c c c@{}}
\toprule
Chapter & Aggregated & Form & Ask & Chat \\
\midrule
Controller \& Contact         & 0.976 (0.005) & 0.975 (0.004) & 0.977 (0.006) & 0.975 (0.005) \\
Purposes \& Legal Bases       & 0.979 (0.008) & 0.979 (0.006) & 0.977 (0.009) & 0.981 (0.008) \\
Personal Data Categories      & 0.973 (0.013) & 0.982 (0.003) & 0.969 (0.015) & 0.970 (0.014) \\
Data Subjects \& Recipients   & 0.971 (0.008) & 0.973 (0.008) & 0.970 (0.008) & 0.971 (0.009) \\
Retention \& Deletion         & 0.949 (0.016) & 0.945 (0.013) & 0.946 (0.018) & 0.955 (0.014) \\
Technical \& Org. Measures    & 0.971 (0.010) & 0.967 (0.012) & 0.972 (0.010) & 0.975 (0.007) \\
\bottomrule
\end{tabular}
\caption{Chapter-level SBERT (ModernBERT) cosine similarity mean
(SD) by mode (full sample $n = 113$).}
\label{tab:nlp-chapter-sbert}
\end{table}

\begin{table}[H]
\centering
\small
\setlength{\tabcolsep}{4pt}
\begin{tabular}{@{}l c c c c@{}}
\toprule
Chapter & Aggregated & Form & Ask & Chat \\
\midrule
Controller \& Contact         & 0.687 (0.122) & 0.663 (0.119) & 0.736 (0.133) & 0.662 (0.101) \\
Purposes \& Legal Bases       & 0.741 (0.074) & 0.763 (0.053) & 0.740 (0.081) & 0.721 (0.080) \\
Personal Data Categories      & 0.713 (0.122) & 0.768 (0.063) & 0.705 (0.149) & 0.669 (0.118) \\
Data Subjects \& Recipients   & 0.706 (0.094) & 0.691 (0.087) & 0.692 (0.115) & 0.733 (0.072) \\
Retention \& Deletion         & 0.519 (0.141) & 0.516 (0.129) & 0.474 (0.150) & 0.566 (0.131) \\
Technical \& Org. Measures    & 0.710 (0.095) & 0.711 (0.073) & 0.701 (0.110) & 0.717 (0.099) \\
\bottomrule
\end{tabular}
\caption{Chapter-level SBERT (\texttt{all-MiniLM-L6-v2}) cosine
similarity mean (SD) by mode (full sample $n = 113$). Reported as
a robustness check on the ModernBERT pass.}
\label{tab:nlp-chapter-sbert-minilm}
\end{table}

\begin{figure}[!htbp]
\centering
\includegraphics[width=0.85\linewidth]{chapter_nlp_bertscore_p.png}\\[1ex]
\includegraphics[width=0.85\linewidth]{chapter_nlp_bertscore_r.png}
\caption{Chapter-level distributions of BERTScore Precision and
Recall by mode (full sample $n = 113$). Continues in
Figure~\ref{fig:chapter-bert-boxplot-b}.}
\label{fig:chapter-bert-boxplot-a}
\end{figure}

\begin{figure}[!htbp]
\centering
\includegraphics[width=0.85\linewidth]{chapter_nlp_bertscore_f1.png}\\[1ex]
\includegraphics[width=0.85\linewidth]{chapter_nlp_sbert.png}
\caption{Chapter-level distributions of BERTScore F1 and SBERT
cosine similarity (ModernBERT) by mode (full sample $n = 113$).
Continued from Figure~\ref{fig:chapter-bert-boxplot-a}.}
\label{fig:chapter-bert-boxplot-b}
\end{figure}
